\chapter[Hipótesis y objetivos]{Hipótesis y objetivos}
\label{cp:hipotesis}
% Extensión estimada: 4--5 páginas.
% Fuentes: PLAN_TFG_GAVIOLA_FACUNDO.pdf (plan aprobado, 25/11/2025: objetivos, metodología en dos fases,
% actividades A--I) y el registro de ejecuciones (Codigo/docs/ejecuciones.md) para fechar las hipótesis que
% surgieron durante el trabajo.
% Criterio: el plan no enuncia hipótesis explícitas; las de partida se derivan de su metodología. Las demás
% se presentan en el orden en que se formularon, para no presentarlas como previstas desde el inicio.

\pendiente{Este capítulo presenta las preguntas de investigación, las hipótesis y los objetivos del trabajo.
Las preguntas y los objetivos provienen del plan de trabajo aprobado. Las hipótesis se separan en dos grupos
según cuándo se formularon: las de partida, implícitas en el plan, y las que surgieron del análisis de los
resultados durante el trabajo. Cada una remite a la sección donde se contrasta.}

\section{Preguntas de investigación}
\label{sec:preguntas}
\pendiente{Las preguntas corresponden a las dos fases de la metodología del plan de trabajo:
\begin{enumerate}[label=\textbf{PI\arabic*.}]
  \item ¿En qué medida un modelo visión--lenguaje aporta mejoras reales frente a modelos orientados
  exclusivamente a visión (CNN y ViT) en el diagnóstico de lesiones de piel, en particular en la
  identificación de melanoma?
  \item ¿Qué tan precisos, completos y clínicamente relevantes resultan los reportes generados
  automáticamente por un VLM ajustado a partir de la imagen?
\end{enumerate}}

\section{Hipótesis}
\label{sec:hipotesis}

\subsection{Hipótesis de partida}
\label{sec:hipotesis_partida}
\pendiente{Estas hipótesis están implícitas en el plan de trabajo, formulado antes de la experimentación:
\begin{enumerate}[label=\textbf{H\arabic*.}]
  \item Un VLM, al integrar información visual y textual, alcanza en el diagnóstico de lesiones de piel un
  desempeño comparable o superior al de una CNN y un ViT entrenados sobre los mismos datos
  (Capítulo~\ref{cp:res_clasif}).
  \item Un VLM ajustado para generación puede producir reportes clínicos precisos y coherentes a partir de la
  imagen (Capítulo~\ref{cp:res_gen}).
\end{enumerate}}

\subsection{Hipótesis surgidas durante el trabajo}
\label{sec:hipotesis_emergentes}
\pendiente{Estas hipótesis se formularon a partir de los resultados y se presentan en el orden en que
surgieron, indicando qué resultado motivó cada una:
\begin{enumerate}[label=\textbf{H\arabic*.}, start=3]
  \item El bajo desempeño de los VLMs generales en clasificación se explica por el desbalance de clases.
  Surge tras refutarse H1.
  \item La especialización de dominio del modelo visión--lenguaje recupera la capacidad de discriminar entre
  clases. Surge tras refutarse H3, al observarse el colapso también en \emph{zero-shot}.
  \item Anclar el generador al diagnóstico de un clasificador especializado (ensamble en cascada) mejora la
  fidelidad diagnóstica de los reportes respecto de la generación directa. Surge al observar que H2 se
  cumple solo parcialmente: los reportes son ricos pero diagnósticamente poco fieles.
  \item Incorporar la incertidumbre del clasificador (afirmar, ofrecer un diferencial o derivar) mejora la
  fidelidad diagnóstica sobre los casos no derivados. Surge del análisis de errores de la cascada.
\end{enumerate}}

\section{Objetivos}
\label{sec:objetivos}

\subsection{Objetivo general}
\pendiente{El objetivo general es desarrollar y evaluar un modelo de
visión--lenguaje capaz de asistir en el diagnóstico de lesiones cutáneas y en la generación automática de
reportes clínicos estructurados a partir de imágenes dermatológicas, con el fin de mejorar la precisión
diagnóstica y optimizar el flujo de trabajo en el estudio de lesiones de piel.}

\subsection{Objetivos específicos}
\pendiente{Los objetivos específicos del plan de trabajo son:
\begin{enumerate}[label=\textbf{OE\arabic*.}]
  \item Diseñar, implementar y evaluar modelos de detección, entre ellos de \emph{deep learning} y VLM, para
  identificar posibles melanomas malignos.
  \item Diseñar y entrenar un modelo visión--lenguaje capaz de integrar información visual y textual para
  producir diagnósticos preliminares y reportes automáticos de lesiones cutáneas.
  \item Evaluar el desempeño del modelo en diferentes escenarios clínicos, analizando su precisión
  diagnóstica, coherencia textual y utilidad práctica para el flujo de trabajo en dermatología.
\end{enumerate}}

\section{Evolución respecto del plan de trabajo}
\label{sec:evolucion_plan}
\pendiente{Esta sección explica, en orden cronológico, cómo se concretó cada objetivo y en qué puntos el
trabajo se apartó del plan, y por qué:
\begin{itemize}
  \item \textbf{OE1.} La comparación CNN--ViT--VLM se amplió a cuatro estrategias de uso de VLMs y, tras su
  colapso, a un modelo visión--lenguaje dermatológico. La detección de melanoma se evalúa como tarea binaria
  derivada de la clasificación en siete clases.
  \item \textbf{OE2.} No se entrena un VLM desde cero: se ajustan VLMs existentes con LoRA sobre datos con
  texto. Como la generación directa resultó diagnósticamente poco fiel, el plan de no comparar la generación
  contra otros enfoques se reemplaza por un ensamble en cascada con el clasificador especializado.
  \item \textbf{OE3.} Los ``escenarios clínicos'' se concretan como robustez a otra población (HIBA) y a
  condiciones de adquisición, calibración de probabilidades y manejo de la incertidumbre. La utilidad
  práctica para el flujo de trabajo no puede establecerse con métricas automáticas.
  \item \textbf{Métricas de generación.} A las métricas de similitud previstas (BERTScore, similitud,
  Jaccard) se agrega la consistencia diagnóstica, porque la similitud textual no detecta reportes con el
  diagnóstico equivocado.
  \item \textbf{Datos.} Se trabaja con HAM10000 (derivado de ISIC 2018) en lugar de PH2, con SkinCaRe y
  MM-Skin para generación y con HIBA para validación externa. Las imágenes histopatológicas quedan fuera
  del alcance.
\end{itemize}}

\section{Alcance y delimitación}
\label{sec:alcance}
\pendiente{El trabajo abarca imágenes dermatoscópicas (HAM10000) como dominio principal, entiende la
detección como clasificación y no como localización, evalúa con métricas automáticas (la validez clínica
requiere especialistas) y se restringe a modelos que corren en GPUs de consumo. No se desarrolla un producto
para uso clínico.}
